\begin{proposition}[6.15.5]
\label{IV.2.s6.15.5-fr}
Sea $f:X\to Y$ un morfismo tal que $f_{\mathrm{red}}$ sea
entero y birracional \textup{(6.15.4)}, y por tanto sobreyectivo.
\begin{enumerate}
\item[\textup{(i)}] Para que $Y$ sea geométricamente unirramoso en un
punto $y$, es necesario y suficiente que $f$ sea radicial en
$y$ y que $X$ sea geométricamente unirramoso en el único
punto $x$ de $f^{-1}(y)$.
\item[\textup{(ii)}] Para que $Y$ sea geométricamente unirramoso, es
necesario y suficiente que $X$ sea geométricamente unirramoso
y que $f$ sea radicial.
\end{enumerate}
\end{proposition}

\begin{proof}
Es posible limitarse, evidentemente, al caso en que $X$ e $Y$
sean \emph{reducidos}; el hecho de que $f$ sea \emph{sobreyectivo} resulta de que $f$
es un morfismo cerrado (\textbf{II}, 6.1.10) y de que $f(X)$ contiene,
por hipótesis, todos los puntos maximales de $Y$. Esto muestra
inmediatamente que (i) implica (ii). Utilizando (\textbf I, 3.6.5)
y (\textbf{II}, 6.1.5), para probar (i) se puede suponer que
$Y=\operatorname{Spec}(\mathcal O_{Y,y})$,
\oldpage[IV-2]{179}
es decir, $Y=\operatorname{Spec}(A)$, donde $A$ es \emph{local}. Como $f$ es
afín, también $X=\operatorname{Spec}(B)$. Si $A$ es
geométricamente unirramoso, es \emph{íntegro}; también lo
es $B$, puesto que $f$ es birracional. Recíprocamente, si $f$
es radicial en $y$ y $X$ es geométricamente unirramoso en $x$,
$B$ tiene un único ideal maximal (pues es entero sobre $A$ y
todo ideal maximal de $B$ se sitúa sobre el de $A$); es decir,
$B$ es un anillo \emph{local}, y decir que $X$ es geométricamente
unirramoso en $x$ significa entonces que $B$ es
geométricamente unirramoso y, por tanto, \emph{íntegro}. Como $f$ es
dominante por hipótesis y $A$ es reducido, se deduce que
$A\subset B$ (\textbf{I}, 1.2.7), de modo que $A$ también es \emph{íntegro}. Así, para probar (i), es posible limitarse
al caso en que $A$ y $B$ son anillos locales y dominios enteros,
$A$ está contenido en $B$ y ambos tienen el mismo cuerpo de
fracciones; por consiguiente, dado que $B$ es entero sobre
$A$, tienen la misma \emph{clausura integral} $C$. La afirmación
resulta entonces de (6.15.3.1, (i)), aplicado a los morfismos
$\operatorname{Spec}(C)\to\operatorname{Spec}(B)$ y
$\operatorname{Spec}(B)\to\operatorname{Spec}(A)$, en virtud de
la interpretación de la propiedad de ser geométricamente
unirramoso en un punto dada en (6.15.3).
\end{proof}
